\chapter{Programs — Statement Code}
%%% generated by the blueprint pipeline from TCSlib.Programs.Light.Compiler.StatementCode
\section{Overview}

Declarations of the Lean module \texttt{TCSlib.Programs.Light.Compiler.StatementCode}.

\section{Declarations}

\begin{definition}[First argument address]
\lean{Light.Compiler.cARG1}
\label{Light.Compiler.cARG1}
\leanok
Define the first argument address to be the integer $-1$.
\end{definition}

\begin{definition}[Second argument address]
\lean{Light.Compiler.cARG2}
\label{Light.Compiler.cARG2}
\leanok
Define the second argument address to be the integer $-2$.
\end{definition}

\begin{definition}[Result address]
\lean{Light.Compiler.cRESULT}
\label{Light.Compiler.cRESULT}
\leanok
Define the result address to be the integer $-3$.
\end{definition}

\begin{definition}[One constant address]
\lean{Light.Compiler.cONE}
\label{Light.Compiler.cONE}
\leanok
Define the address for the constant one to be the integer $-5$.
\end{definition}

\begin{definition}[Zero constant address]
\lean{Light.Compiler.cZERO}
\label{Light.Compiler.cZERO}
\leanok
Define the address for the constant zero to be the integer $-7$.
\end{definition}

\begin{definition}[Negation address]
\lean{Light.Compiler.cNEG}
\label{Light.Compiler.cNEG}
\leanok
Define the address for negation to be the integer $-9$.
\end{definition}

\begin{definition}[Frame pointer offset]
\lean{Light.Compiler.cFP}
\label{Light.Compiler.cFP}
\leanok
Define the frame-pointer offset to be the integer $-11$.
\end{definition}

\begin{definition}[Address offset]
\lean{Light.Compiler.cADDR}
\label{Light.Compiler.cADDR}
\leanok
Define the address offset to be the integer $-13$.
\end{definition}

\begin{definition}[Return offset]
\lean{Light.Compiler.cRET}
\label{Light.Compiler.cRET}
\leanok
Define the return offset to be the integer $-15$.
\end{definition}

\begin{definition}[Return-register offset]
\lean{Light.Compiler.cRR}
\label{Light.Compiler.cRR}
\leanok
Define the return-register offset to be the integer $-17$.
\end{definition}

\begin{definition}[Next frame pointer offset]
\lean{Light.Compiler.cNFP}
\label{Light.Compiler.cNFP}
\leanok
Define the next frame-pointer offset to be the integer $-19$.
\end{definition}

\begin{definition}[Displacement offset]
\lean{Light.Compiler.cD}
\label{Light.Compiler.cD}
\leanok
Define the displacement offset to be the integer $-21$.
\end{definition}

\begin{definition}[Indexed offset]
\lean{Light.Compiler.cT}
\label{Light.Compiler.cT}
\leanok
For each natural number $k$, define the indexed offset to be $-(23+4k)$ as an integer.
\end{definition}

\begin{definition}[Pool offset]
\lean{Light.Compiler.cPool}
\label{Light.Compiler.cPool}
\leanok
For each natural number $n$, define the pool offset to be $-(25+4n)$ as an integer.
\end{definition}

\begin{definition}[Stack register addresses]
\lean{Light.Compiler.cStack}
\label{Light.Compiler.cStack}
\leanok
For each natural number $j$, define the corresponding stack address to be the integer
$-2j$.
\end{definition}

\begin{definition}[Main frame size]
\lean{Light.Compiler.mainFrame}
\label{Light.Compiler.mainFrame}
\leanok
The size of the main frame is $4$.
\end{definition}

\begin{definition}[Instruction for an arithmetic operation]
\lean{Light.Op.instr}
\label{Light.Op.instr}
\leanok
\uses{Light.Op, WordRAM.Instr}
An arithmetic operation is translated to the WordRAM instruction with the same
operation: addition to addition, subtraction to subtraction, and multiplication to
multiplication. The instruction takes the three supplied integer register addresses as
its destination and two source addresses, respectively.
\end{definition}

\begin{definition}[Code for a digit list]
\lean{Light.Compiler.digitsCode}
\label{Light.Compiler.digitsCode}
\leanok
\uses{Light.Compiler.cONE, WordRAM.Instr}
Fix an integer register address $c$. For a finite list of natural numbers, define its
code recursively: the empty list produces the instruction that subtracts register $c$
from itself and stores the result in $c$; for a list with first entry $d$ and remaining
list $L$, append to the code for $L$ the instruction that doubles register $c$, and
then, if $d$ is nonzero, the instruction that adds the contents of register $-5$ to
register $c$.
\end{definition}

\begin{definition}[Length assigned to a digit list]
\lean{Light.Compiler.digitsLen}
\label{Light.Compiler.digitsLen}
\leanok
Assign a natural number to each finite list of natural numbers as follows: the empty
list has length $1$; a list with first entry $d$ and tail $L$ has assigned length equal
to the assigned length of $L$, plus $1$, plus another $1$ when $d$ is nonzero.
\end{definition}

\begin{definition}[Compilation of expressions]
\lean{Light.Compiler.compileExpr}
\label{Light.Compiler.compileExpr}
\leanok
\uses{Light.Compiler.cADDR, Light.Compiler.cFP, Light.Compiler.cPool, Light.Compiler.cT, Light.Compiler.digitsCode, Light.Expr, Light.Op.instr, WordRAM.Instr}
For an expression and a natural-number temporary-register index $k$, compile a constant
$n$ by applying the digit-list code at address $-(23+4k)$ to the base-$2$ digits of $n$.
Compile a variable $x$ by subtracting the pool address $-(25+8x)$ from the frame-pointer
register address $-11$ and storing the result at address $-13$, then loading into
address $-(23+4k)$ from the memory address held there. For an operation applied to
expressions $a$ and $b$, concatenate their compilations at indices $k$ and $k+1$, then
apply the operation to the registers at addresses $-(23+4k)$ and $-(23+4(k+1))$, storing
the result at the first address; for a load applied to $a$, append to its compilation at
index $k$ the instruction that loads into that register from the memory address held in
it.
\end{definition}

\begin{definition}[Size of an expression]
\lean{Light.Expr.size}
\label{Light.Expr.size}
\leanok
\uses{Light.Compiler.digitsLen, Light.Expr}
Define the size of an expression recursively: a constant has size equal to the assigned
length of its base-$2$ digit list; a variable has size $2$; an operation applied to $a$
and $b$ has size $\text{size}(a)+\text{size}(b)+1$; and a load applied to $a$ has size
$\text{size}(a)+1$.
\end{definition}

\begin{definition}[Compilation of conditions]
\lean{Light.Compiler.compileCond}
\label{Light.Compiler.compileCond}
\leanok
\uses{Light.Compiler.cD, Light.Compiler.cONE, Light.Compiler.cT, Light.Compiler.cZERO, Light.Compiler.compileExpr, Light.Cond, WordRAM.Instr}
For a condition and a natural-number jump position $l$, compile a comparison of $a$ and
$b$ by first compiling $a$ at temporary index $0$ and $b$ at index $1$. For the
condition $a<b$, subtract the value in register $0$ from the value in register $1$,
subtract the contents of register $-5$, and jump to position $l$ if the result is
negative. For the condition $a=b$, subtract the value in register $1$ from the value in
register $0$ and jump to $l$ if negative; then replace the result by the contents of
register $-7$ minus that result and jump to $l$ if negative.
\end{definition}

\begin{definition}[Condition size]
\lean{Light.Cond.size}
\label{Light.Cond.size}
\leanok
\uses{Light.Cond, Light.Expr.size}
For an expression, define its size recursively: a constant has size equal to the length
of its binary digit list, where the empty list has length $1$ and each digit contributes
$1$ plus another $1$ if it is nonzero; a variable has size $2$, a binary operation has
size the sum of its operands’ sizes plus $1$, and a load has size its operand’s size
plus $1$. The size of a less-than condition is the sum of its two expressions’ sizes
plus $3$, while the size of an equality condition is that sum plus $4$.
\end{definition}

\begin{definition}[Procedure code layout]
\lean{Light.Compiler.CodeLayout}
\label{Light.Compiler.CodeLayout}
\leanok
A code layout consists of a natural number $F$ giving the number of local variables in
each frame, a list of natural-number procedure positions, and a natural-number position
for the dispatcher.
\end{definition}

\begin{definition}[Argument code generation]
\lean{Light.Compiler.argsCode}
\label{Light.Compiler.argsCode}
\leanok
\uses{Light.Compiler.compileExpr, Light.Expr, WordRAM.Instr}
Given a list of expressions and a starting temporary-slot index $k$, generate code by
concatenating the code for each expression in order, compiling the first at slot $k$ and
each successive expression at the next slot. The empty list generates no instructions.
\end{definition}

\begin{definition}[Total expression size]
\lean{Light.Compiler.argsSize}
\label{Light.Compiler.argsSize}
\leanok
\uses{Light.Expr, Light.Expr.size}
For a list of expressions, its total size is the sum of the sizes of its expressions;
the empty list has size $0$.
\end{definition}

\begin{definition}[Frame initialization code]
\lean{Light.Compiler.frameInit}
\label{Light.Compiler.frameInit}
\leanok
\uses{Light.Compiler.cADDR, Light.Compiler.cNFP, Light.Compiler.cPool, Light.Compiler.cT, Light.Compiler.cZERO, WordRAM.Instr}
Given a frame size $F$ and an input length $n$, generate instructions to initialize the
$F$ local-variable cells: cell $i$ receives input value $i$ when $i<n$, and otherwise
receives zero, for each $0\le i<F$.
\end{definition}

\begin{definition}[New frame setup code]
\lean{Light.Compiler.newFrame}
\label{Light.Compiler.newFrame}
\leanok
\uses{Light.Compiler.cADDR, Light.Compiler.cFP, Light.Compiler.cNFP, Light.Compiler.cPool, Light.Compiler.cZERO, Light.Compiler.frameInit, WordRAM.Instr}
Given a frame size $F$, an input length $n$, and a return position, generate
instructions that allocate a new frame, initialize its $F$ local-variable cells from the
first $n$ input values and fill the remaining cells with zero, store the return position
in the frame, and set the frame pointer to the new frame.
\end{definition}

\begin{definition}[Procedure call code]
\lean{Light.Compiler.callCode}
\label{Light.Compiler.callCode}
\leanok
\uses{Light.Compiler.CodeLayout, Light.Compiler.argsCode, Light.Compiler.cNEG, Light.Compiler.newFrame, Light.Expr, WordRAM.Instr}
Given a code layout, a procedure index $p$, a list of arguments, and a return position,
generate the argument code starting at temporary slot $0$, then set up a new frame of
the layout’s frame size with those arguments and the given return position, and finally
branch to the procedure position at index $p$ in the layout’s entry list. If that index
is out of range, branch to position $0$.
\end{definition}

\begin{definition}[Call code size]
\lean{Light.Compiler.callSize}
\label{Light.Compiler.callSize}
\leanok
\uses{Light.Compiler.argsSize, Light.Expr}
For a frame size $F$ and a list of arguments, the call size is the total size of the
arguments plus $1+2F+4$.
\end{definition}

\begin{definition}[Statement size measure]
\lean{Light.Stmt.size}
\label{Light.Stmt.size}
\leanok
\uses{Light.Compiler.callSize, Light.Cond.size, Light.Expr.size, Light.Stmt}
For a natural number $F$, define the size of a statement recursively as follows. The
skip statement has size $0$; an assignment has size the expression size plus $2$; a
store has size the sum of its address and value expression sizes plus $1$; a sequence
has size the sum of its parts; a conditional has size the condition size plus the two
branch sizes plus $1$; a while loop has size the condition size plus the body size plus
$1$; and a procedure call has size the sum of its argument expression sizes, plus
$2F+7$. Here expression size is $2$ for a variable, the length measure of the binary
digits for a constant (where the empty list has length measure $1$ and each digit
contributes $1$ plus $1$ if it is nonzero), the sum of the two operand sizes plus $1$
for a binary operation, and the operand size plus $1$ for a load; condition size is the
sum of its two expression sizes plus $3$ for a strict comparison and plus $4$ for
equality.
\end{definition}

\begin{definition}[Statement compilation]
\lean{Light.Compiler.compileStmt}
\label{Light.Compiler.compileStmt}
\leanok
\uses{Light.Compiler.CodeLayout, Light.Compiler.cADDR, Light.Compiler.cFP, Light.Compiler.cNEG, Light.Compiler.cPool, Light.Compiler.cRET, Light.Compiler.cT, Light.Compiler.callCode, Light.Compiler.callSize, Light.Compiler.compileCond, Light.Compiler.compileExpr, Light.Cond.size, Light.Stmt, Light.Stmt.size, WordRAM.Instr}
Given a code layout with $F$ local variables per frame, a list of procedure entry
positions, and a dispatcher position, the compiler maps a statement and a starting code
position to a list of Word RAM instructions. It emits no instructions for skip;
evaluates and stores expressions for assignments and stores; concatenates compiled code
for sequences, placing the second statement after the first according to their sizes;
and emits condition checks, branch code, and jumps for conditionals and loops. For a
call, it evaluates the arguments, creates and initializes a new frame, jumps to the
selected procedure entry (or position $0$ if the procedure index is out of range), then
stores the returned value in the designated variable.
\end{definition}

\begin{definition}[Procedure epilogue]
\lean{Light.Compiler.epilogue}
\label{Light.Compiler.epilogue}
\leanok
\uses{Light.Compiler.CodeLayout, Light.Compiler.cADDR, Light.Compiler.cFP, Light.Compiler.cNEG, Light.Compiler.cPool, Light.Compiler.cRET, Light.Compiler.cRR, WordRAM.Instr}
Given a code layout with $F$ local variables per frame and dispatcher position $d$, the
epilogue is the five-instruction sequence that loads the return address from the current
frame, loads the return value from the caller's frame, restores the caller's frame
pointer by moving it back by the frame size, and branches to position $d$.
\end{definition}

\begin{definition}[Epilogue instruction count]
\lean{Light.Compiler.epilogueSize}
\label{Light.Compiler.epilogueSize}
\leanok
The epilogue size is $5$.
\end{definition}

\begin{definition}[Procedure dispatcher]
\lean{Light.Compiler.dispatcher}
\label{Light.Compiler.dispatcher}
\leanok
\uses{Light.Compiler.cONE, Light.Compiler.cRR, WordRAM.Instr}
For each integer $j$ with $0 \langle j<n$, the dispatcher appends two instructions:
decrement the return register by one, then branch to code position $j$ if that register
is negative. Thus the dispatcher is the concatenation of these pairs in increasing order
of $j$.
\end{definition}

\begin{theorem}[Length of Digit Code]
\lean{Light.Compiler.length_digitsCode}
\label{Light.Compiler.length_digitsCode}
\leanok
\uses{Light.Compiler.digitsCode, Light.Compiler.digitsLen}
\difficulty{4}
For an integer register $c$ and a finite list $L$ of natural-number digits, the digit
code has length $D(L)$, where $D([])=1$ and $D(d::L)=D(L)+1$ if $d=0$, while
$D(d::L)=D(L)+2$ if $d\ne0$.
\end{theorem}

\begin{theorem}[Length of Compiled Expressions]
\lean{Light.Compiler.length_compileExpr}
\label{Light.Compiler.length_compileExpr}
\leanok
\uses{Light.Compiler.compileExpr, Light.Compiler.length_digitsCode, Light.Expr, Light.Expr.size}
\difficulty{4}
For every expression $e$ and temporary-register index $k\in\mathbb N$, the compiled
expression has length equal to its size. The size is the binary-digit code length for a
constant, $2$ for a variable, the sum of the operand sizes plus $1$ for a binary
arithmetic operation, and the operand size plus $1$ for a memory load.
\end{theorem}

\begin{theorem}[Length of Compiled Conditions]
\lean{Light.Compiler.length_compileCond}
\label{Light.Compiler.length_compileCond}
\leanok
\uses{Light.Compiler.compileCond, Light.Compiler.length_compileExpr, Light.Cond, Light.Cond.size}
\difficulty{4}
For every condition $c$ and branch target $\ell\in\mathbb N$, the compiled condition has
length equal to the condition's size. For a comparison $a<b$ this size is
$\mathrm{size}(a)+\mathrm{size}(b)+3$, and for an equality test $a=b$ it is
$\mathrm{size}(a)+\mathrm{size}(b)+4$.
\end{theorem}

\begin{theorem}[Argument code length]
\lean{Light.Compiler.length_argsCode}
\label{Light.Compiler.length_argsCode}
\leanok
\uses{Light.Compiler.argsCode, Light.Compiler.argsSize, Light.Compiler.length_compileExpr, Light.Expr}
\difficulty{4}
For any finite list $es$ of expressions and any natural number $k$, the instruction
sequence generated for the arguments starting at index $k$ has length equal to the sum
of the sizes of the expressions in $es$.
\end{theorem}

\begin{theorem}[Frame initialization length]
\lean{Light.Compiler.length_frameInit}
\label{Light.Compiler.length_frameInit}
\leanok
\uses{Light.Compiler.frameInit}
\difficulty{1}
For all natural numbers $F$ and $n$, the frame initialization list has length $2F$.
\end{theorem}

\begin{theorem}[New frame length]
\lean{Light.Compiler.length_newFrame}
\label{Light.Compiler.length_newFrame}
\leanok
\uses{Light.Compiler.length_frameInit, Light.Compiler.newFrame}
\difficulty{3}
For all natural numbers $F$, $n$, and $r$, the new frame list with $F$ local slots, $n$
initialized values, and return position $r$ has length $2F+4$.
\end{theorem}

\begin{theorem}[Length of Compiled Calls]
\lean{Light.Compiler.length_callCode}
\label{Light.Compiler.length_callCode}
\leanok
\uses{Light.Compiler.CodeLayout, Light.Compiler.callCode, Light.Compiler.callSize, Light.Compiler.length_argsCode, Light.Compiler.length_newFrame, Light.Expr}
\difficulty{3}
Let $G$ be a code layout, $p\in\mathbb N$ a procedure index, $args$ a finite list of
expressions, and $r\in\mathbb N$ a return position. The compiled call, which compiles
the arguments, constructs a new frame of width $G.F$, and branches to the indexed
procedure entry (or position $0$ if no such entry exists), has length $\sum_{e\in
args}\mathrm{size}(e)+1+2G.F+4$.
\end{theorem}

\begin{theorem}[Length of the Procedure Epilogue]
\lean{Light.Compiler.length_epilogue}
\label{Light.Compiler.length_epilogue}
\leanok
\uses{Light.Compiler.CodeLayout, Light.Compiler.epilogue, Light.Compiler.epilogueSize}
\difficulty{1}
For every code layout $G$, the procedure epilogue has length $5$ instructions.
\end{theorem}

\begin{theorem}[Length of Compiled Statements]
\lean{Light.Compiler.length_compileStmt}
\label{Light.Compiler.length_compileStmt}
\leanok
\uses{Light.Compiler.CodeLayout, Light.Compiler.compileStmt, Light.Compiler.length_callCode, Light.Compiler.length_compileCond, Light.Compiler.length_compileExpr, Light.Stmt, Light.Stmt.size}
\difficulty{4}
For every code layout $G$, statement $s$, and starting position $p\in\mathbb N$, the
compiled code for $s$ starting at $p$ has length equal to the size of $s$ for frame
width $G.F$.
\end{theorem}

\begin{definition}[Variables occurring in an expression]
\lean{Light.Expr.vars}
\label{Light.Expr.vars}
\leanok
\uses{Light.Expr}
For an expression, define its variable bound recursively: a constant has bound $0$, a
variable with index $x$ has bound $x+1$, a binary operation has bound the maximum of the
bounds of its two operands, and a load has the bound of its address expression.
\end{definition}

\begin{definition}[Height of an expression]
\lean{Light.Expr.height}
\label{Light.Expr.height}
\leanok
\uses{Light.Expr}
For expressions built from natural-number constants and variables using binary addition,
subtraction, and multiplication, and unary loads, define the height $h$ recursively by
$h(c)=h(x)=1$ for a constant $c$ or variable $x$, $h(\text{op}(a,b))=\max(h(a),h(b)+1)$
for a binary operation with operands $a,b$, and $h(\text{load}(a))=h(a)$.
\end{definition}

\begin{definition}[Expression width]
\lean{Light.Expr.width}
\label{Light.Expr.width}
\leanok
\uses{Light.Expr, Light.Expr.height, Light.Expr.vars}
For an expression $e$, its width is the maximum of its variable bound and its height.
\end{definition}

\begin{definition}[Width of a condition]
\lean{Light.Cond.width}
\label{Light.Cond.width}
\leanok
\uses{Light.Cond, Light.Expr.width}
The width of a less-than condition or an equality condition is the maximum of the widths
of its two expressions.
\end{definition}

\begin{definition}[Maximum over a list]
\lean{Light.maxOf}
\label{Light.maxOf}
\leanok
Given a function from a type to the natural numbers, define its maximum over a list
recursively: the empty list has maximum $0$, and the maximum for a nonempty list is the
larger of the function value at the first element and the maximum over the remaining
list.
\end{definition}

\begin{definition}[Width of a statement]
\lean{Light.Stmt.width}
\label{Light.Stmt.width}
\leanok
\uses{Light.Cond.width, Light.Expr.width, Light.Stmt, Light.maxOf}
The width of a statement is defined recursively. It is $0$ for a skip; for an assignment
it is the maximum of the assigned variable's index plus one and the expression width;
for a store it is the maximum of the address and value expression widths; for a sequence
it is the maximum of the two statement widths; for a conditional it is the maximum of
the condition width and both branch widths; and for a loop it is the maximum of the
condition and body widths. For a procedure call, it is the maximum of the destination
variable's index plus one, the number of arguments, and the maximum width of the
argument expressions (with maximum $0$ when there are no arguments).
\end{definition}

\begin{theorem}[Bound on a list maximum]
\lean{Light.maxOf_le_iff}
\label{Light.maxOf_le_iff}
\leanok
\uses{Light.maxOf}
\difficulty{4}
Let $f$ map a type $\alpha$ to the natural numbers, let $l$ be a list of elements of
$\alpha$, and let $F$ be a natural number. The recursively defined maximum of the values
of $f$ on $l$ (with value $0$ for the empty list) is at most $F$ if and only if $f(y)\le
F$ for every $y$ in $l$.
\end{theorem}

\begin{theorem}[Variable bound for a binary operation]
\lean{Light.Expr.vars_op_le_iff}
\label{Light.Expr.vars_op_le_iff}
\leanok
\uses{Light.Expr.vars}
\difficulty{1}
For any binary operation symbol $o$ and expressions $a,b$, the variable bound of the
expression formed by applying $o$ to $a$ and $b$ is at most $F$ if and only if the
variable bounds of both $a$ and $b$ are at most $F$.
\end{theorem}

\begin{theorem}[Expression width bound]
\lean{Light.Expr.width_le_iff}
\label{Light.Expr.width_le_iff}
\leanok
\uses{Light.Expr.height, Light.Expr.vars, Light.Expr.width}
\difficulty{1}
For every expression $e$ and natural number $F$, the width of $e$ is at most $F$ if and
only if both its variable bound and its height are at most $F$.
\end{theorem}

\begin{theorem}[Width bound for strict inequality conditions]
\lean{Light.Cond.width_lt_le_iff}
\label{Light.Cond.width_lt_le_iff}
\leanok
\uses{Light.Cond.width, Light.Expr.width}
\difficulty{1}
For expressions $a,b$ and a natural number $F$, the width of the strict-inequality
condition comparing $a$ and $b$ is at most $F$ if and only if both $a$ and $b$ have
width at most $F$.
\end{theorem}

\begin{theorem}[Width bound for equality conditions]
\lean{Light.Cond.width_eq_le_iff}
\label{Light.Cond.width_eq_le_iff}
\leanok
\uses{Light.Cond.width, Light.Expr.width}
\difficulty{1}
For expressions $a,b$ and a natural number $F$, the width of the equality condition
comparing $a$ and $b$ is at most $F$ if and only if both $a$ and $b$ have width at most
$F$.
\end{theorem}

\begin{theorem}[Width bound for assignment]
\lean{Light.Stmt.width_set_le_iff}
\label{Light.Stmt.width_set_le_iff}
\leanok
\uses{Light.Expr.width, Light.Stmt.width}
\difficulty{1}
For an assignment to variable $x$ with expression $e$, its width is at most $F$ if and
only if $x<F$ and $e$ has width at most $F$.
\end{theorem}

\begin{theorem}[Width bound for memory store]
\lean{Light.Stmt.width_store_le_iff}
\label{Light.Stmt.width_store_le_iff}
\leanok
\uses{Light.Expr.width, Light.Stmt.width}
\difficulty{1}
For a memory store with address expression $a$ and value expression $e$, its width is at
most $F$ if and only if $a$ and $e$ each have width at most $F$.
\end{theorem}

\begin{theorem}[Width bound for sequential composition]
\lean{Light.Stmt.width_seq_le_iff}
\label{Light.Stmt.width_seq_le_iff}
\leanok
\uses{Light.Stmt.width}
\difficulty{1}
For the sequential composition of statements $s$ and $t$, its width is at most $F$ if
and only if each of $s$ and $t$ has width at most $F$.
\end{theorem}

\begin{theorem}[Width bound for a conditional]
\lean{Light.Stmt.width_ite_le_iff}
\label{Light.Stmt.width_ite_le_iff}
\leanok
\uses{Light.Cond.width, Light.Stmt.width}
\difficulty{1}
For a conditional statement with condition $c$ and branches $s$ and $t$, its width is at
most $F$ if and only if the widths of $c$, $s$, and $t$ are each at most $F$.
\end{theorem}

\begin{theorem}[Width bound for a while loop]
\lean{Light.Stmt.width_while_le_iff}
\label{Light.Stmt.width_while_le_iff}
\leanok
\uses{Light.Cond.width, Light.Stmt.width}
\difficulty{1}
For a while loop with condition $c$ and body $s$, its width is at most $F$ if and only
if the widths of $c$ and $s$ are each at most $F$.
\end{theorem}

\begin{theorem}[Width bound for a procedure call]
\lean{Light.Stmt.width_call_le_iff}
\label{Light.Stmt.width_call_le_iff}
\leanok
\uses{Light.Expr.width, Light.Stmt.width, Light.maxOf_le_iff}
\difficulty{1}
For a procedure call with target variable $x$ and argument list $A$, its width is at
most $F$ if and only if $x<F$, the number of arguments is at most $F$, and every
argument has width at most $F$.
\end{theorem}
